% ==============================================================================
% Fichier : 06_post_exploitation.tex
% Description : Post-Exploitation, Escalade de Privilèges et Mouvement Latéral
% ==============================================================================

\chapter{Post-Exploitation et Mouvement Latéral}

Obtenir un shell initial (ex: utilisateur "www-data") n'est souvent que le début. La phase de \textbf{Post-Exploitation} vise à consolider l'accès, augmenter les droits (Escalade) et se déplacer vers des cibles plus sensibles (Mouvement Latéral).

\section{Escalade de Privilèges (Privilege Escalation)}
% ------------------------------------------------------------------------------

L'escalade de privilèges est l'art de passer d'un utilisateur standard à un utilisateur avec des droits d'administration (Root sous Linux, SYSTEM/Admin sous Windows).

\subsection{Linux : De l'Utilisateur à Root}

\begin{itemize}
    \item \textbf{Kernel Exploits :} Exploitation de failles dans le noyau de l'OS (ex: Dirty COW, Dirty Pipe). Risque de crash système.
    \item \textbf{SUID/SGID Binaries :} Fichiers exécutables qui tournent avec les droits de leur propriétaire (souvent Root). Si un éditeur de texte (vim, less) a le bit SUID, on peut lire le fichier \texttt{/etc/shadow}.
    \item \textbf{Sudo Misconfiguration :} Droit d'exécuter des commandes précises en root sans mot de passe (ex: \texttt{sudo vim} permet de lancer un shell).
    \item \textbf{Cron Jobs :} Tâches planifiées exécutées par Root. Si un script appelé est modifiable par l'utilisateur, il peut injecter du code.
\end{itemize}

\begin{lstlisting}[language=bash, caption=Trouver les fichiers SUID]
find / -perm -u=s -type f 2>/dev/null
\end{lstlisting}

\subsection{Windows : De l'Utilisateur à SYSTEM}

\begin{itemize}
    \item \textbf{Unquoted Service Path :} Si le chemin d'un service contient des espaces et n'est pas entre guillemets (ex: \texttt{C:\Program Files\My App\service.exe}), Windows cherche l'exécutable dans l'ordre. En plaçant un maliciel à \texttt{C:\Program.exe}, on peut le faire exécuter par SYSTEM.
    \item \textbf{AlwaysInstallElevated :} Si cette politique est activée, n'importe quel utilisateur peut installer un package MSI (.msi) avec les droits SYSTEM.
    \item \textbf{Token Impersonation :} Techniques comme "Rotten Potato" ou "PrintSpoofer" permettant de voler le jeton d'un processus système.
\end{itemize}

\begin{boxlizbiz}
\textbf{Scénario LiZbiZ :}
Vous avez obtenu un shell "apache" sur le serveur Web Linux. Vous listez les permissions sudo : \texttt{sudo -l}. Résultat : l'utilisateur apache peut lancer \texttt{/usr/bin/vim}.
\textbf{Exploitation :} Ouvrir vim en sudo, puis exécuter \texttt{:!bash} dans l'éditeur. Vous êtes maintenant "root".
\end{boxlizbiz}

% ------------------------------------------------------------------------------
\section{Récupération d'Identifiants (Credential Harvesting)}
% ------------------------------------------------------------------------------

L'objectif est de voler des mots de passe pour se connecter ailleurs (mot de passe souvent réutilisé).

\subsection{Dumping de Mémoire (Mimikatz)}

Sous Windows, les mots de passe sont stockés en mémoire (LSASS). \textbf{Mimikatz} est l'outil roi pour les extraire.

\begin{lstlisting}[language=bash, caption=Commandes Mimikatz basiques]
privilege::debug
sekurlsa::logonpasswords
# Affiche les mots de passe en clair ou les hashs NTLM des utilisateurs connectes
\end{lstlisting}

\subsection{Fichiers de Configuration}

Les applications stockent souvent les mots de passe en clair :
\begin{itemize}
    \item Web : \texttt{web.config}, \texttt{wp-config.php} (WordPress), \texttt{database.yml}.
    \item SSH : Clés privées dans \texttt{~/.ssh/id\_rsa}.
\end{itemize}

% ------------------------------------------------------------------------------
\section{Mouvement Latéral (Lateral Movement)}
% ------------------------------------------------------------------------------

Se déplacer de machine en machine sans repasser par l'extérieur.

\subsection{Pass-the-Hash (PtH)}

Technique permettant de s'authentifier avec le hash NTLM du mot de passe au lieu du mot de passe en clair. Très efficace si le dump Mimikatz a fonctionné.

\begin{lstlisting}[language=bash, caption=Pass-the-Hash avec CrackMapExec]
crackmapexec smb 192.168.1.0/24 -u Administrator -H "A9FD...HASH..." 
# Teste le hash sur tout le reseau pour trouver ou il est valide
\end{lstlisting}

\subsection{Outils d'Exécution à Distance}

\begin{itemize}
    \item \textbf{PsExec :} Outil Microsoft légitime souvent détecté par les AV.
    \item \textbf{WMI / WinRM :} Protocoles de gestion Windows natifs, très furtifs.
    \item \textbf{Impacket :} Suite Python d'outils réseaux (ex: \texttt{wmiexec.py}, \texttt{psexec.py}).
\end{itemize}

% ------------------------------------------------------------------------------
\section{Attaques Active Directory}
% ------------------------------------------------------------------------------

L'Active Directory est le "Saint Graal" de l'entreprise. Il centralise les identités.

\subsection{BloodHound}

Outil de cartographie des relations AD. Il visualise graphiquement "Qui peut administrer quoi".
\textbf{Utilité :} Trouver le chemin le plus court vers le compte "Domain Admin".

\subsection{Kerberoasting}

Attaque qui consiste à demander un ticket de service (TGS) pour un compte de service, puis à le cracker hors-ligne (offline) pour récupérer le mot de passe du compte de service.

% ------------------------------------------------------------------------------
\section{Pivotage (Pivoting)}
% ------------------------------------------------------------------------------

Comment accéder à un réseau interne (VLAN 2) depuis une machine compromise en DMZ (VLAN 1) ?

\begin{boxdef}
\textbf{Pivoting :} Utiliser une machine compromise comme relais (proxy) pour scanner et attaquer des réseaux inaccessibles depuis la machine de l'attaquant.
\end{boxdef}

\begin{lstlisting}[language=bash, caption=Configuration Proxychains]
# Dans /etc/proxychains.conf
[ProxyList]
socks4 127.0.0.1 1080  # La connexion passe par le tunnel SSH/Empire etabli sur la cible
\end{lstlisting}

\begin{lstlisting}[language=bash, caption=Scan via le pivot]
proxychains nmap -sT -Pn 10.0.0.5
# Scan un reseau interne (10.0.0.x) via la machine compromise
\end{lstlisting}

\begin{boxlizbiz}
\textbf{Scenario Final LiZbiZ :}
1. Compromission du serveur Web (DMZ).
2. Pivotage vers le réseau interne via l'outil \texttt{Chisel} ou \texttt{Metasploit autoroute}.
3. Scan du réseau interne : découverte du Contrôleur de Domaine (DC).
4. Utilisation de Mimikatz sur un poste Windows compromis pour dumper les identifiants.
5. Pass-the-Hash vers le DC : \textbf{Contrôle total de l'entreprise LiZbiZ}.
\end{boxlizbiz}
